\documentclass[11pt]{article}
\usepackage[margin=1in]{geometry}
\usepackage{amsmath,amssymb,amsthm}
\usepackage{graphicx}
\usepackage[hidelinks]{hyperref}

\newtheorem{theorem}{Theorem}
\newtheorem{lemma}{Lemma}
\newtheorem{remark}{Remark}

\newcommand{\R}{\mathbb R}
\newcommand{\A}{\mathcal A}
\newcommand{\G}{\mathcal G}
\newcommand{\norm}[1]{\left\lVert #1\right\rVert}
\newcommand{\eps}{\varepsilon}
\newcommand{\dd}{\,\mathrm d}

\title{Two-Sided Inverse-Gaussian Poisson Square Bounds Characterize UMD}
\author{FA Banach run \texttt{fa\_banach\_001}, agent \texttt{agent\_lane\_10}}
\date{23 July 2026}

\begin{document}
\maketitle

\section*{Status}

This packet records a candidate full solution, likely valid, to the question
in the remark following Theorem 1.5 of Almeida, Betancor, Fari\~na, and
Rodr\'iguez-Mesa \cite{ABFR}. The answer is affirmative.

\section*{Source Passage}

\begin{center}
\includegraphics[width=0.98\linewidth]{figures/open_problem_crop.png}
\end{center}

\section*{Transcription and Logical Scope of the Question}

Let \(X\) be an order-continuous K\"othe function space, let \(1<p<\infty\),
and write
\[
 \G_{\A}^X f(x,w)
 =
 \left(\int_0^\infty
 \left|t\partial_tP_t^\A(f(\,\cdot\,,w))(x)\right|^2
 \frac{\dd t}{t}\right)^{1/2}.
\]
Theorem 1.5(ii) of \cite{ABFR} assumes an upper estimate for
\(\G_\A^X\) on \(L_X^p(\R^n,\gamma_{-1})\) and an upper estimate for the
dual-space square function on \(L_{X^*}^{p'}\). The source asks whether the
second, dual-space estimate can instead be replaced by
\[
 \norm{f}_{L_X^p(\R^n,\gamma_{-1})}
 \le C\norm{\G_\A^X f}_{L_X^p(\R^n,\gamma_{-1})},
 \tag{R}\label{eq:reverse}
\]
and whether UMD then follows.

Thus the modified condition retains the first estimate in Theorem 1.5(ii):
\[
 \norm{\G_\A^X f}_{L_X^p(\R^n,\gamma_{-1})}
 \le C\norm{f}_{L_X^p(\R^n,\gamma_{-1})}.
 \tag{U}\label{eq:upper}
\]
The question is about the two-sided pair \eqref{eq:upper}--\eqref{eq:reverse},
not about \eqref{eq:reverse} alone.

\section*{Idea of the Proof}

The inverse-Gaussian geometry becomes Euclidean at small scales. Test
\eqref{eq:upper}--\eqref{eq:reverse} on a shrinking bump
\(f_\eps(x)=F(x/\eps)\). After the changes \(x=\eps z\) and \(t=\eps s\),
the generator becomes
\[
 B_\eps=-\tfrac12\Delta-\eps^2z\cdot\nabla,
\]
the measure becomes \(e^{\eps^2|z|^2}\dd z\), and \(B_\eps\) tends to the
Euclidean generator \(-\Delta/2\).

The required convergence of the full square-function norms is not left to an
abstract semigroup limit. On a fixed physical ball, the local comparison
estimate proved in \cite{ABFR} gives a rescaled error of size
\(\eps^{1/2}I_{1/2}(\norm{F}_X)\), whose norm tends to zero for every
\(p>1\). Outside that ball, the input and output are separated. The explicit
global kernel estimate (A2) in \cite{ABFR} supplies an
\(L^p(\gamma_{-1})\) Gaussian majorant, so the normalized tail is
\(O(\eps^{n(1-1/p)})\).

Consequently the two inverse-Gaussian inequalities pass to two-sided
Euclidean Poisson square-function inequalities. The Xu--Hyt\"onen
characterization then gives UMD.

\section*{The Tangent Lemma}

Let \(P_t\) denote the Euclidean Poisson semigroup subordinated to the heat
semigroup generated by \(-\Delta/2\), and put
\[
 \G_0^X F(z,w)
 =
 \left(\int_0^\infty
 |t\partial_tP_t(F(\,\cdot\,,w))(z)|^2\frac{\dd t}{t}\right)^{1/2}.
\]
The harmless choice of the factor \(1/2\) does not change the square
function: a constant rescaling of \(t\) leaves \(\dd t/t\) invariant.

\begin{lemma}[Small-scale tangent lemma]\label{lem:tangent}
Let \(1<p<\infty\), let \(X\) be a K\"othe function space, and let
\[
 F(z,w)=\sum_{j=1}^N\phi_j(z)x_j(w),
 \qquad \phi_j\in C_c^\infty(\R^n),\quad x_j\in X.
\]
For \(0<\eps<1\), define \(f_\eps(x,w)=F(x/\eps,w)\). If
\(\dd\gamma_{-1}(x)=c_n e^{|x|^2}\dd x\), then
\[
 \lim_{\eps\downarrow0}
 \eps^{-n/p}\norm{f_\eps}_{L_X^p(\gamma_{-1})}
 =c_n^{1/p}\norm{F}_{L_X^p(\dd z)}
 \tag{1}\label{eq:input-limit}
\]
and
\[
 \lim_{\eps\downarrow0}
 \eps^{-n/p}\norm{\G_\A^Xf_\eps}_{L_X^p(\gamma_{-1})}
 =c_n^{1/p}\norm{\G_0^XF}_{L_X^p(\dd z)}.
 \tag{2}\label{eq:square-limit}
\]
\end{lemma}

\begin{proof}
The first limit follows immediately from
\[
 \eps^{-n}\norm{f_\eps}_{L_X^p(\gamma_{-1})}^p
 =
 c_n\int_{\R^n}\norm{F(z)}_X^p e^{\eps^2|z|^2}\dd z
\]
and dominated convergence, since \(F\) is compactly supported.

We prove \eqref{eq:square-limit}. Choose \(R>0\) with
\(\operatorname{supp}F\subset B(0,R)\), and then choose
\[
 0<\rho<\min\{1,n/4\}.
\]
For all sufficiently small \(\eps\), if \(|x|<\rho\) and
\(y\in\eps B(0,R)\), then \((x,y)\) belongs to the local region \(N_1\)
of \cite{ABFR}. The local comparison estimate proved there is
\[
 \left\|t\partial_t\big(P_t^\A(x,y)-P_t(x-y)\big)
 \right\|_{L^2((0,\infty),\dd t/t)}
 \le
 C\frac{\sqrt{1+|x|}}{|x-y|^{\,n-1/2}}.
 \tag{3}\label{eq:local-source}
\]
The reverse triangle inequality in \(L^2(\dd t/t)\), followed by Minkowski's
inequality in the lattice, therefore gives, for \(|z|<\rho/\eps\),
\begin{align*}
 &\norm{\G_\A^Xf_\eps(\eps z)-\G_0^XF(z)}_X\\
 &\quad\le
 C\eps^{1/2}\int_{B(0,R)}
 \frac{\norm{F(u)}_X}{|z-u|^{\,n-1/2}}\dd u
 =:C\eps^{1/2}J(z).
 \tag{4}\label{eq:local-error}
\end{align*}
Here we also used the exact Euclidean scaling
\(\G_0^Xf_\eps(\eps z)=\G_0^XF(z)\).

The function \(J\) is bounded on bounded sets because
\(|v|^{-n+1/2}\) is locally integrable, while
\[
 J(z)\le C_F|z|^{-n+1/2}\qquad (|z|>2R).
\]
Moreover, \(e^{\eps^2|z|^2}\le e^{\rho^2}\) on the ball under
consideration. Hence the \(p\)-th power of the norm in
\eqref{eq:local-error}, with the rescaled weight, is bounded by a constant
times
\[
 \eps^{p/2}\left(
 1+\int_1^{\rho/\eps}
 r^{\,n-1-p(n-1/2)}\dd r\right).
 \tag{5}\label{eq:radial}
\]
If \(p(n-1/2)>n\), this is \(O(\eps^{p/2})\). At equality it is
\(O(\eps^{p/2}\log(1/\eps))\). If \(p(n-1/2)<n\), it is
\(O(\eps^{n(p-1)})\). Thus \eqref{eq:local-error} tends to zero in
\(L_X^p(B(0,\rho/\eps),e^{\eps^2|z|^2}\dd z)\) for every \(p>1\).

It remains to show that the inverse-Gaussian output outside the fixed
physical ball is negligible after normalization. Define the scalar square
kernel
\[
 K_\A(x,y)=
 \left\|t\partial_tP_t^\A(x,y)\right\|_{L^2((0,\infty),\dd t/t)}.
\]
We claim that, after decreasing the allowed size of \(y\) if necessary, there
is a function \(q\in L^p(\R^n,\gamma_{-1})\) such that
\[
 K_\A(x,y)\le q(x),\qquad |x|\ge\rho,\quad |y|\le\rho/2.
 \tag{6}\label{eq:majorant}
\]
On a fixed annulus \(\rho\le|x|\le R_0\), this follows from
\eqref{eq:local-source} together with the Euclidean bound
\[
 \left\|t\partial_tP_t(x-y)\right\|_{L^2(\dd t/t)}
 \le C|x-y|^{-n}
\]
on the local part, and from the source global bound on the complementary
part. All pairs are uniformly separated there.

For \(R_0\) sufficiently large, every pair with
\(|x|>R_0\) and \(|y|\le\rho/2\) is global. The Poisson-to-heat comparison
in \cite{ABFR}, followed by its estimate (A2) with \(m=1,k=0,q=2\), gives
parameters \(\theta,\delta>0\), chosen so that
\(\theta+\delta>2/p\), and the two bounds
\[
 K_\A(x,y)\le
 C e^{-(\theta+\delta)|x|^2/2+(\theta-\delta)|y|^2/2}
 \quad\text{if }\langle x,y\rangle\le0,
 \tag{7}\label{eq:A2negative}
\]
and
\[
 K_\A(x,y)\le
 C|x+y|^n
 e^{\theta(|y|^2-|x|^2)/2
   -\delta|x+y||x-y|/2}
 \quad\text{if }\langle x,y\rangle>0.
 \tag{8}\label{eq:A2positive}
\]
Since
\[
 |x+y||x-y|\ge\big||x|^2-|y|^2\big|,
\]
both cases, uniformly for bounded \(y\), are dominated for large \(x\) by
\[
 C(1+|x|)^n e^{-a|x|^2}
\quad\text{for some }a>1/p.
\]
This belongs to \(L^p(\gamma_{-1})\), proving
\eqref{eq:majorant}.

Let \(M_F=\int_{B(0,R)}\norm{F(u)}_X\dd u\). Minkowski and
\eqref{eq:majorant} give, for small \(\eps\) and \(|x|\ge\rho\),
\[
 \norm{\G_\A^Xf_\eps(x)}_X
 \le \int_{\eps B(0,R)}K_\A(x,y)
       \norm{F(y/\eps)}_X\dd y
 \le \eps^nM_Fq(x).
\]
Consequently,
\[
 \eps^{-n/p}
 \norm{\G_\A^Xf_\eps}_{L_X^p(\{|x|\ge\rho\},\gamma_{-1})}
 \le
 M_F\norm q_{L^p(\gamma_{-1})}\eps^{n(1-1/p)}
 \longrightarrow0.
 \tag{9}\label{eq:tail}
\]

Finally, \(\G_0^XF\in L_X^p(\R^n)\): it is dominated by a finite sum of
scalar Poisson square functions times \(\norm{x_j}_X\). On
\(|z|<\rho/\eps\), the weight is bounded by \(e^{\rho^2}\), converges
pointwise to one, and the domains increase to \(\R^n\). Dominated convergence
therefore yields
\[
 \int_{|z|<\rho/\eps}\norm{\G_0^XF(z)}_X^p
 e^{\eps^2|z|^2}\dd z
 \longrightarrow
 \norm{\G_0^XF}_{L_X^p(\R^n)}^p.
\]
Combine this with the local error estimate and \eqref{eq:tail}, then undo
the change \(x=\eps z\), to obtain \eqref{eq:square-limit}.
\end{proof}

\section*{Full Answer}

\begin{theorem}\label{thm:answer}
Let \(X\) be an order-continuous K\"othe function space and
\(1<p<\infty\). The following are equivalent.
\begin{enumerate}
\item \(X\) is UMD.
\item There is a constant \(C\) such that, for every
\(f\in L_X^p(\R^n,\gamma_{-1})\),
\[
 C^{-1}\norm f_{L_X^p(\gamma_{-1})}
 \le\norm{\G_\A^Xf}_{L_X^p(\gamma_{-1})}
 \le C\norm f_{L_X^p(\gamma_{-1})}.
\]
\end{enumerate}
In particular, the proposed replacement of the \(X^*\)-condition in
Theorem 1.5(ii) of \cite{ABFR} is valid.
\end{theorem}

\begin{proof}
Assume the two inverse-Gaussian estimates. Apply them to the functions
\(f_\eps\) in Lemma \ref{lem:tangent}, divide by \(\eps^{n/p}\), and let
\(\eps\downarrow0\). Equations \eqref{eq:input-limit} and
\eqref{eq:square-limit} give, on compactly supported finite tensors,
\[
 C^{-1}\norm F_{L_X^p(\R^n)}
 \le\norm{\G_0^XF}_{L_X^p(\R^n)}
 \le C\norm F_{L_X^p(\R^n)}.
 \tag{10}\label{eq:euclidean-two-sided}
\]
The upper estimate and the reverse triangle inequality for the
\(L^2(\dd t/t)\) norm extend the square-function map continuously from this
dense test class to all of \(L_X^p(\R^n)\); the lower estimate passes to the
completion as well.

Xu's Banach-lattice theorem \cite[Theorem 4.1]{Xu}, equivalently the
Banach-lattice specialization of Hyt\"onen's characterization
\cite[Theorem 1.6]{Hytonen}, says that the two-sided Euclidean Poisson
lattice-square estimate \eqref{eq:euclidean-two-sided} implies that \(X\) is
UMD. This proves the direction asked in the source remark.

Conversely, if \(X\) is UMD, then the two-sided inverse-Gaussian estimate is
the UMD implication for subordinated symmetric diffusion semigroups recalled
as Theorem D(i) in \cite{ABFR}. Hence the conditions are equivalent.
\end{proof}

\section*{Verification Notes}

The proof uses only the local comparison estimate, the Poisson-to-heat
comparison, and the explicit global estimate (A2) already proved in the
source paper, followed by the classical Euclidean two-sided-square
characterization of UMD.

The local radial integral \eqref{eq:radial} was checked separately in all
three exponent regimes. The far-tail exponent is strict because the
parameters in (A2) are selected in the source with
\(\theta+\delta>2/p\). The order-continuity assumption is exactly the source
assumption; no stronger lattice hypothesis is introduced.

No numerical experiment is used. A separate line-by-line audit is included
as \texttt{verification.md}.

\section*{Novelty and Literature Check}

The lightweight run indexes were searched for arXiv:2102.13381, the paper
title, inverse-Gaussian square functions, reverse estimates, UMD, and removal
of the dual-space hypothesis. No existing result packet answers this exact
question.

Bounded external searches on 23 July 2026 used the exact phrase ``replace the
condition involving the dual space'', the full title together with ``UMD
reverse estimate'', and close inverse-Gaussian square-function
formulations. They returned the source paper and its Potential Analysis
publication, but no later paper explicitly resolving the remark. The
open-access published version was also inspected directly: the question
remains verbatim on published PDF page 43, so it was not removed or answered
during peer review. This supports, but cannot guarantee, novelty.

\section*{Limitations and Human-Review Recommendation}

The theorem answers the source question exactly. It does not assert that the
reverse estimate \eqref{eq:reverse} alone implies UMD; the retained upper
estimate \eqref{eq:upper} is essential to the statement being solved.

Human review should focus on Lemma \ref{lem:tangent}, especially the
uniform majorant \eqref{eq:majorant} extracted from (A2). The packet spells
out both sign cases of (A2), the required strict exponent, the compact
off-diagonal region, and the normalization of the shrinking tail.

\begin{thebibliography}{9}

\bibitem{ABFR}
V. Almeida, J. J. Betancor, J. C. Fari\~na, and L. Rodr\'iguez-Mesa,
\emph{Littlewood-Paley-Stein theory and Banach spaces in the inverse Gaussian
setting}, Potential Analysis 59 (2023), 1235--1284;
arXiv:2102.13381.

\bibitem{Hytonen}
T. P. Hyt\"onen,
\emph{Littlewood-Paley-Stein theory for semigroups in UMD spaces},
Rev. Mat. Iberoam. 23 (2007), no. 3, 973--1009.

\bibitem{Xu}
Q. Xu,
\emph{Littlewood-Paley theory for functions with values in uniformly convex
spaces}, J. Reine Angew. Math. 504 (1998), 195--226.

\end{thebibliography}

\end{document}
